\section{L'IA comme outil : applications}

% ---------- Analyse de données ----------
\domaincard{Analyser les données de la neuroscience}{
    \begin{itemize}
        \item Automatiser des tâches très longues à la main
        \item Cartographier des cerveaux, trier des signaux, mesurer le comportement
    \end{itemize}
}{
    \begin{itemize}
        \item Des outils déjà utilisés dans de nombreux laboratoires
        \item Exemple : reconstruction du câblage complet d'un cerveau
    \end{itemize}
}

\begin{frame}{Cartographier un cerveau entier : FlyWire}
\begin{itemize}
    \item Le plan de câblage complet d'un cerveau de mouche adulte : \textbf{139\,255 neurones} et \textbf{54,5 millions de synapses}
    \item Construit à partir d'images de microscopie électronique, avec l'IA pour reconstruire les neurones et des humains pour les corriger
    \item Une ressource ouverte pour formuler des hypothèses testables
\end{itemize}
% IMAGE : rendu 3D des neurones de la mouche (à ajouter, avec crédit)
\src{Dorkenwald et al. (FlyWire), Nature 2024}
\end{frame}

% ---------- Décodage ----------
\domaincard{Décoder le langage depuis le cerveau}{
    \begin{itemize}
        \item Redonner la parole à des personnes paralysées
        \item Contrôler un ordinateur ou une prothèse par la pensée
    \end{itemize}
}{
    \begin{itemize}
        \item Des essais cliniques avec des électrodes implantées (voir l'ouverture)
        \item Des démonstrations non invasives, encore limitées au laboratoire
    \end{itemize}
}

\begin{frame}{Décoder le sens sans chirurgie : fMRI}
\begin{itemize}
    \item Un décodeur reconstruit le \textbf{sens} d'une histoire entendue ou imaginée, sans reproduire les mots exacts
    \item Il faut \textbf{16 heures} d'histoires par participant pour entraîner le décodeur
    \item Le décodeur s'appuie sur un modèle de langage et sur un modèle d'encodage cérébral
    \item La coopération est nécessaire : le participant peut le faire échouer en résistant consciemment
\end{itemize}
\src{Tang, LeBel, Jain \& Huth, Nature Neuroscience 2023}
\end{frame}

\begin{frame}{Décoder avec des capteurs externes : Brain2Qwerty}
\begin{itemize}
    \item 35 volontaires sains tapent des phrases sur un clavier pendant un enregistrement
    \item Erreur sur les caractères : \textbf{32\%} avec la MEG, 67\% avec l'EEG, et 19\% pour les meilleurs participants
    \item Limites : décode des gestes de frappe, impossibles pour des patients paralysés ; MEG non portable ; calibrage long et propre à chaque personne
\end{itemize}
\src{Lévy et al. (Meta), 2025}
\end{frame}

% ---------- Clinique ----------
\domaincard{Clinique : aider au diagnostic}{
    \begin{itemize}
        \item Repérer tôt des marqueurs de langage de maladies neurologiques
        \item Classer les aphasies et suivre la récupération
    \end{itemize}
}{
    \begin{itemize}
        \item Des études de recherche sur des paroles de patients
        \item Pas encore un outil clinique de routine : jeux de données petits, généralisation limitée
    \end{itemize}
}

\begin{frame}{Classer l'aphasie progressive primaire}
\begin{itemize}
    \item Des modèles de langage analysent la parole de patients (78 personnes)
    \item Ils séparent trois variantes avec \textbf{88,5\%} d'accord avec les diagnostics cliniques indépendants
    \item Les groupes trouvés correspondent à des régions d'atrophie du cortex
\end{itemize}
\src{Rezaii et al., préprint medRxiv 2023}
\end{frame}

\begin{frame}{La surprise d'un modèle comme marqueur}
\begin{itemize}
    \item La surprise calculée par des modèles de langage détecte la présence d'aphasie et ses différents types
    \item Données : AphasiaBank, une base d'enregistrements de personnes aphasiques
    \item Les modèles complètent les indices de langage déjà utilisés par les cliniciens
\end{itemize}
\src{Cong et al., Scientific Reports 2024}
\end{frame}

\begin{frame}{Un enjeu d'équité : quelles langues ?}
\begin{itemize}
    \item Les troubles du langage varient selon la langue du patient
    \item L'anglais est un mauvais étalon pour les autres langues
    \item Les langues africaines sont presque absentes des données et des outils
\end{itemize}
\begin{block}{Piste}
    Constituer des données de parole et de langage dans nos langues, pour ne pas importer des diagnostics calibrés sur d'autres.
\end{block}
% TODO : ajouter les langues de Guinée (pular, malinké, soussou) si tu veux ancrer l'exemple
\src{Revue dans Brain 2023, 146(12)}
\end{frame}

% ---------- Jumeaux numériques ----------
\domaincard{Jumeaux numériques et contrôle}{
    \begin{itemize}
        \item Un modèle assez fidèle remplace des expériences coûteuses
        \item Choisir les stimuli qui activent ou éteignent une région
    \end{itemize}
}{
    \begin{itemize}
        \item Un jumeau numérique du cortex visuel de la souris
        \item Des phrases choisies pour piloter le réseau du langage humain
    \end{itemize}
}

\begin{frame}{Un jumeau numérique du cortex visuel}
\begin{itemize}
    \item Modèle entraîné sur l'activité de souris qui regardent des films
    \item Il prédit les réponses de dizaines de milliers de neurones à de \textbf{nouvelles vidéos}
    \item Il prédit même des propriétés anatomiques, validées sur des images de microscopie électronique de la même souris
    \item Limite : seulement le cortex visuel de la souris, en observation passive
\end{itemize}
\src{Wang et al. (MICrONS), Nature 2025}
\end{frame}

\begin{frame}{Piloter le réseau du langage}
\begin{itemize}
    \item Un modèle de type GPT ajusté sur les réponses fMRI à 1\,000 phrases
    \item Il sert à \textbf{choisir de nouvelles phrases} qui augmentent ou suppriment l'activité du réseau du langage
    \item L'expérience confirme les prédictions chez de nouveaux participants
\end{itemize}
\src{Tuckute et al., Nature Human Behaviour 2024}
\end{frame}
